\documentclass[a4paper,5pt]{amsbook}
%%%%%%%%%%%%%%%%%%%%%%%%%%%%%%%%%%%%%%%%%%%%%%%%%%%%%%%%%%%%%%%%%%%%%

\usepackage{booktabs}
\usepackage{graphicx}
\usepackage{multicol}
\usepackage{textcomp}
\usepackage{systeme}
\usepackage{amssymb}
\usepackage[]{amsmath}
\usepackage{subcaption}
\usepackage[inline]{enumitem}
\usepackage{gensymb}
\usepackage[utf8]{inputenc}

%%%%%%%%%%%%%%%%%%%%%%%%%%%%%%%%%%%%%%%%%%%%%%%%%%%%%%%%%%%%%%

\renewcommand{\sin}{\,\mbox{sen}\,}
\newcommand{\sen}{\,\mbox{sen}\,}
\newcommand{\tg}{\,\mbox{tg}\,}
\newcommand{\cosec}{\,\mbox{cosec}\,}
\newcommand{\cotg}{\,\mbox{cotg}\,}
\newcommand{\tr}{\,\mbox{tr}\,}
\newcommand{\ds}{\displaystyle}

%%%%%%%%%%%%%%%%%%%%%%%%%%%%%%%%%%%%%%%%%%%%%%%%%%%%%%%%%%%%%%%%%%%%%%%%

\setlength{\textwidth}{16cm} \setlength{\topmargin}{-1.3cm}
\setlength{\textheight}{30cm}
\setlength{\leftmargin}{1.2cm} \setlength{\rightmargin}{1.2cm}
\setlength{\oddsidemargin}{0cm}\setlength{\evensidemargin}{0cm}

%%%%%%%%%%%%%%%%%%%%%%%%%%%%%%%%%%%%%%%%%%%%%%%%%%%%%%%%%%%%%%%%%%%%%%%%

% \renewcommand{\baselinestretch}{1.6}
% \renewcommand{\thefootnote}{\fnsymbol{footnote}}
% \renewcommand{\theequation}{\thesection.\arabic{equation}}
% \setlength{\voffset}{-50pt}
% \numberwithin{equation}{chapter}

%%%%%%%%%%%%%%%%%%%%%%%%%%%%%%%%%%%%%%%%%%%%%%%%%%%%%%%%%%%%%%%%%%%%%%%

\begin{document}
\thispagestyle{empty}
\pagestyle{empty}
\begin{minipage}[h]{0.14\textwidth}
    \includegraphics[scale=0.24]{../../ufgd.png}
\end{minipage}
\begin{minipage}[h]{\textwidth}
    \begin{tabular}{c}
        {{\bf UNIVERSIDADE FEDERAL DA GRANDE DOURADOS}}\\
        {{\bf C\'{a}lculo Diferencial e Integral II --- Lista 9}}\\
        {{\bf Prof.\ Adriano Barbosa}}\\
    \end{tabular}
    \vspace{-0.45cm}
    %
\end{minipage}

%------------------------

\vspace{1cm}
%%%%%%%%%%%%%%%%%%%%%%%%%%%%%%%%   formulario  inicio  %%%%%%%%%%%%%%%%%%%%%%%%%%%%%%%%
\begin{enumerate}
\vspace{0.5cm}
\item Determine se as s\'eries s\~ao convergentes ou divergentes
	\begin{enumerate}
        \vspace{0.3cm}
		\item $\ds\sum_{n=1}^\infty \frac{{(2n+1)}^n}{n^{2n}}$
        \vspace{0.3cm}
		\item $\ds\sum_{k=1}^\infty \frac{1}{k^\pi}$
        \vspace{0.3cm}
		\item $\ds\sum_{k=1}^\infty k^2 e^{-k}$
        \vspace{0.3cm}
		\item $\ds\sum_{n=1}^\infty \frac{n!}{e^{n^2}}$
	\end{enumerate}

\vspace{0.5cm}
\item Encontre o raio e o intervalo de converg\^encia das s\'eries
	\begin{enumerate}
        \vspace{0.3cm}
		\item $\ds\sum_{n=0}^\infty \frac{n{(x+2)}^n}{3^{n+1}}$
        \vspace{0.3cm}
		\item $\ds\sum_{n=1}^\infty \frac{{(x-2)}^n}{n^n}$
	\end{enumerate}

\vspace{0.5cm}
\item
	\begin{enumerate}
        \vspace{0.3cm}
        \item Escreva as fun\c{c}\~oes $\sin x$ e $\cos x$ como s\'erie de Maclaurin e
            encontre seu raio e intervalo de converg\^encia.
        \vspace{0.3cm}
        \item Utilize o item $(a)$ e a s\'erie de Maclaurin da fun\c{c}\~ao $e^x$ para
            verificar a F\'ormula de Euler: $e^{ix} = \cos x + i\ \sin x$, onde $i$
            \'e a unidade imagin\'aria.
	\end{enumerate}

\vspace{0.5cm}
\item Encontre a s\'erie de Taylor das fun\c{c}\~oes abaixo centradas no valor dado
	\begin{enumerate}
        \vspace{0.3cm}
		\item $f(x) = x^4 - 3x^2 + 1$, $a = 1$
        \vspace{0.3cm}
		\item $f(x) = \ln x$, $a = 2$
	\end{enumerate}

\end{enumerate}
\end{document}
